% Generated draft: requires mathematical and source review.
\documentclass[11pt]{article}
\usepackage[margin=25mm]{geometry}
\usepackage{fontspec}
\setmainfont{Latin Modern Roman}
\usepackage{amsmath,amssymb,mathtools,unicode-math}
\setmathfont{Latin Modern Math}
\usepackage{enumitem,needspace,xurl,hyperref}
\setlength{\parindent}{0pt}
\setlength{\parskip}{5pt}
\setlength{\emergencystretch}{3em}
\newcommand{\N}{\mathbb{N}}
\newcommand{\Z}{\mathbb{Z}}
\newcommand{\Q}{\mathbb{Q}}
\newcommand{\R}{\mathbb{R}}
\newcommand{\C}{\mathbb{C}}
\newcommand{\F}{\mathbb{F}}
\newcommand{\A}{\mathbb{A}}
\newcommand{\PP}{\mathbb{P}}
\newcommand{\E}{\mathbb{E}}
\newcommand{\eps}{\varepsilon}
\newcommand{\dd}{\,\mathrm{d}}
\begin{document}
\section*{UnsolvedMath ID: 801123}
\textbf{Source problem code:} catalogue entry
\subsection{Definitions and mathematical statement}
A type \(\mathrm{II}_1\) factor is a von Neumann algebra \(M\) with center \(Z(M)=\C 1\), a faithful normal tracial state \(\tau\) satisfying \(\tau(1)=1\), and no nonzero minimal projections. Property \((T)\) is understood in its von Neumann algebra sense: there exist a finite set \(F\subset M\) and \(\varepsilon>0\) such that every Hilbert \(M\)-\(M\) bimodule with commuting normal unital left and right actions and a unit vector \(\xi\) satisfying
\[
\|x\xi-\xi x\|<\varepsilon,\qquad
\bigl|\langle x\xi,\xi\rangle-\tau(x)\bigr|<\varepsilon,\qquad
\bigl|\langle \xi x,\xi\rangle-\tau(x)\bigr|<\varepsilon
\quad(x\in F)
\]
contains a nonzero central vector, meaning a vector \(\eta\ne0\) with \(x\eta=\eta x\) for every \(x\in M\).

The supplied question is: does there exist a type \(\mathrm{II}_1\) factor \(M\) with property \((T)\) and a unique sequential commutation orbit?

\textbf{Unresolved qualification.} The supplied record does not define ``sequential commutation orbit.'' In particular, it does not identify the objects forming the orbits, the relation or action defining them, or any required ambient algebra or ultrapower. Consequently, the final condition cannot be made into a precise mathematical predicate from this record alone. No separability or ultrapower hypothesis has been added; any such intended requirements must be recovered from the cited source before this draft is mathematically complete.
\subsection{Short English statement}
Elayavalli's Problem 24 asks whether a property \((T)\) type \(\mathrm{II}_1\) factor can have a unique sequential commutation orbit. The supplied record leaves that orbit notion undefined, so the existence question requires further source clarification. Its ``open'' status is a supplied registry label.
\subsection{Sources}
\begin{itemize}
\item Local UnsolvedMath registry, supplied complete source record: permanent ID \(801123\), source problem code \texttt{catalogue entry}, title ``Ultraproduct \(\mathrm{II}_1\) Factors --- Problem 24.'' The record attributes the problem to Elayavalli, with proposed year 2025, and labels it open. Upstream revision: \texttt{d371ee917ac9647fbb06cd6fc529deec6f3dda82}.
\item Supplied source list: Elayavalli, \emph{50 Open Problems: Ultraproduct \(\mathrm{II}_1\) Factors} (2025), Problem 24. Supplied URL: \url{https://arxiv.org/abs/2511.20377}. The record reports access on 2026-10-16 and extraction by \texttt{arxiv\_latex}. These are supplied provenance claims; the source and its mathematical status have not been independently verified here.
\end{itemize}
\end{document}
